
\begin{ledgroupsized}[r]{120mm}
\footnotesize
\pstart
\noindent\textbf{Überlieferung:}
\pend
\end{ledgroupsized}
\begin{ledgroupsized}[r]{114mm}
\footnotesize
\pstart \parindent -6mm
\makebox[6mm][l]{\textit{L}}%
Konzept:
LH~XXXV~10, 14~Bl.~2. 
Ein Blatt~4\textsuperscript{o}
\href{\Scan{LH_35_10_14_002wz}}%
{Wasserzeichenfragment} am Blattrand;
alle Ränder beschnitten.
Eineinhalb Seiten.
In der oberen linken Ecke von Bl.~2~r\textsuperscript{o}, in Rot, der Vermerk von unklarer Hand \textit{:\;II}
(vgl. die Träger der benachbarten Signaturen LH~XXXV~10, 14 Bl.~2 bis 13, u.\,a. 
N.\;\ref{60117},
\ref{60118},
\ref{60115},
\ref{60112},
\ref{60116},
\ref{60113},
\ref{60114},
\ref{60110}
und \ref{60111}
in diesem Band).
\pend
\end{ledgroupsized}
\selectlanguage{latin}
\vspace{8mm}
\pstart
\centering
\href{\Scan{LH_35_10_14_002r}}%
{\lbrack2~r\textsuperscript{o}\rbrack}
\hfill
\edtext{%
\edtext{Ex Relatione}{%
\lemma{}%
\Afootnote{%
\textit{Oberhalb der Überschrift nachgetragen:} nicht abgangen}}
\edtext{mea}{%
\lemma{}%
\Bfootnote{%
mea 
\textit{erg.~L}}}
de novissimo \cite{02001}libro \indexN{\textls{Vanni}, Giovanni Francesco 1638\textendash1709}P.\,Vanni 
de momentis gravium 
\hfill
\phantom{\lbrack2~r\textsuperscript{o}\rbrack}
\linebreak
1694 \cite{01023}\textit{Actis} Lips. inserenda}{%
\lemma{Ex relatione \lbrack...\rbrack\ inserenda}%
\Cfootnote{Gemeint ist Leibnizens Rezension von \indexN{\textls{Vanni}, Giovanni Francesco 1638\textendash1709}\textsc{G.\,F. Vanni}, \cite{02001}\textit{Investigatio momentorum quibus gravia tendunt deorsum}, Rom 1693, N.\;\ref{61037} in diesem Band.}}
\pend
\vspace{0.5\baselineskip}
\pstart
\noindent 
\edtext{Argumentum in his \textit{Actis} oppositum	%
\edtext{\indexN{\textls{Vanni}, Giovanni Francesco 1638\textendash1709}\lbrack R.\,P.\;Vannio\rbrack}{%
\lemma{\lbrack R.\,P.\;Vannio\rbrack}%
\Cfootnote{Eckige Klammern von Leibniz.}}
quo etiam \indexN{\textls{Torricelli} (Torricellius), Evangelista 1608\textendash1647}Torricellius jam usus hic refertur,
quod gravia\indexS{gravia in aequilibrio} tunc censenda sint esse in aequilibrio\indexS{aequilibrium}	%
cum commune eorum centrum gravitatis\indexS{centrum gravitatis commune gravium} 
descendere non potest,\indexS{centrum gravitatis descendere nequiens}
eo effugere satagit \indexN{\textls{Vanni}, Giovanni Francesco 1638\textendash1709}noster 		%
quod duo gravia, filo connexa,\indexS{gravia duo filo connexa} et in duobus diversis planis descendere\indexS{gravia in diversis planis descendere conantia} 		%
\edtext{conantia communi}{%
\lemma{conantia}%
\Bfootnote{%
\textit{(1)}~(ut ipse arb
\textit{(2)}~communi~\textit{L}}}
gravitatis centro,\indexS{centrum gravitatis commune} (ut ipse arbitratur) careant.%
\indexS{gravia centro gravitatis commune carentia}}{%
\lemma{Argumentum \lbrack...\rbrack\ careant}%
\Cfootnote{Die Passage stimmt mit dem letzten Absatz der gedruckten Rezension weitgehend wörtlich überein (\refnote{FnC:AE_1694_112-113_argumentum}).
Leibniz spielt auf seine anonym erschienene \textit{Demonstratio geometrica regulae apud Staticos receptae} an (\textit{AE} November 1685, S.\;501\textendash505, N.\;\ref{49409} in diesem Band, bes. \refpassage{AE_1685_501-505_15a}{AE_1685_501-505_15b}.
Allerdings erwähnt \indexN{\textls{Vanni}, Giovanni Francesco 1638\textendash1709}Vanni 
in diesem Zusammenhang keine Publikation aus den \cite{01023}\textit{Acta Eruditorum}, sondern \indexN{\textls{Torricelli} (Torricellius), Evangelista 1608\textendash1647}Torricellis \textit{De motu gravium}: siehe die Erläuterung zu N.\;\ref{61037}, \refnote{FnC:AE_1694_112-113_Torricelli}.}}
\pend
\pstart
Haec sequentia addere constitueram, sed postea omisi, ne controversia relationi misceretur.
\pend
\pstart
\edtext{\indexN{\textls{Leibniz} (Leibnitius, GGL), Gottfried Wilhelm 1646\textendash1716}Autor autem objectionis, intellecta hac responsione nobis significavit}{%
\lemma{Autor \lbrack...\rbrack\ significavit}%
\Cfootnote{Der anonyme Rezensent (\indexN{\textls{Leibniz} (Leibnitius, GGL), Gottfried Wilhelm 1646\textendash1716}Leibniz) stellt die nachfolgenden Ausführungen nicht als die eigenen dar,
sondern referiert angeblich aus einer Mitteilung des ebenfalls anonymen Autors der \textit{Demonstratio geometrica}, der wiederum \indexN{\textls{Leibniz} (Leibnitius, GGL), Gottfried Wilhelm 1646\textendash1716}Leibniz selbst ist.}}
\textls{centrum gravitatis commune}\indexS{centrum gravitatis commune plurium corporum}
\edtext{plurium corporum}{%
\lemma{}%
\Bfootnote{plurium corporum \textit{erg.~L}}}
a se hic intelligi quod mente assignatur quovis instanti,
perinde ac si corpora, vel eorum centra gravitatis\indexS{centra gravitatis propria corporum}
\edtext{\textls{propria}\lbrack,\rbrack\ essent conjuncta 
lineis gravitate carentibus,\indexS{linea gravitate carens} iisque rigidis,\indexS{linea rigida} 
licet de momento in momentum (ob privatos corporum motus)\indexS{motus privati corporum} 
variantibus sive aliter assignandis; atque}{%
\lemma{\textls{propria}}%
\Bfootnote{%
\textit{(1)}~lineis rigidis, licet de momento in momentum variantibus, sed gravitate carentibus
\textit{(2)}~gravit
\textit{(3)}~lineis \lbrack...\rbrack\ momentum
\textit{(a)}~variantibus
\textit{(b)}~(ob
\textit{(aa)}~privatus
\textit{(bb)}~privatos \lbrack...\rbrack\ assignandis essent conjuncta;
\textit{(4)}~essent conjuncta \lbrack...\rbrack\ assignandis; atque~\textit{L}}}
ita hoc sensu nihil referre utrum corpora gravia
\edtext{sint connexa}{%
\lemma{}%
\Bfootnote{%
sint connexa 
\textit{erg.~L}}}
firmiter,\indexS{gravia firmiter connexa} vel ut hic 
\edtext{flexiliter;\indexS{gravia flexiliter connexa} an}{%
\lemma{flexiliter;}%
\Bfootnote{%
\textit{(1)}~vel
\textit{(2)}~an~\textit{L}}}
vero sint omnino dissoluta.\indexS{gravia omnino dissoluta} 
Quoniam utique semper rectae centra gravium\indexS{centrum gravitatis} duorum conjungentes,\indexS{recta centra duorum gravium conjungens}
mente ductae, intelligi, et secundum legem aequilibrii\indexS{leges aequilibrii} secari possunt, quo facto punctum sectionis\indexS{punctum sectionis}
dicetur centrum commune gravitatis;\indexS{centrum gravitatis commune} 
etsi id non sit affixum corporibus, sed fixum ad ea variet,
quando ipsa sibi rigide connexa\indexS{gravia non rigide connexa} non sunt.
Isque loquendi usus jam apud eruditos\indexS{eruditi} est receptus.
Qua definitione posita cum negari nequeat
tale punctum\indexS{punctum sectionis} assignari posset, 
negabit Cl.~\indexN{\textls{Vanni}, Giovanni Francesco 1638\textendash1709}Vannius suppositionem,
quod scilicet tali centro\indexS{centrum gravitatis} vel puncto\indexS{punctum sectionis} 
descendere nequeunte\indexS{centrum gravitatis descendere nequiens} 
sit aequilibrium.\indexS{aequilibrium}
Replicatur ergo ab \indexN{\textls{Leibniz} (Leibnitius, GGL), Gottfried Wilhelm 1646\textendash1716}autore
objectionis non tantum nullam in contrarium instantiam repertum iri,
sed etiam posse suppositionem hanc demonstrari
ex eo, quod ubicunque hoc centrum gravitatis totale\indexS{centrum gravitatis totale} descendit,
ibi semper reperitur plus de gravibus\indexS{gravia plus descendentia quam ascendentia} descendere, quam ascendit, 
et quando hoc punctum descendere nequit,\indexS{centrum gravitatis descendere nequiens} 
tunc vel nihil, vel non plus descendere,
quam quantum ascendit, adeoque nihil agi.
Et revera considerationem
\href{\Scan{LH_35_10_14_002v}}%
{\lbrack2~v\textsuperscript{o}\rbrack}
centri gravitatis\indexS{consideratio centri gravitatis} 
tantum pro compendio hujus considerationis videri ab antiquis\indexS{antiqui} introductam.
Itaque missa omni mentione centri
\edtext{gravitatis,\indexS{centrum gravitatis} \textls{receptam}}{%
\lemma{gravitatis,}%
\Bfootnote{%
\textit{(1)}~statuit objector
\textit{(2)}~\textls{receptam}~\textit{L}}}
\textls{legem aequilibrii in planis}
\edtext{\textls{inclinatis\indexS{lex aequilibrii in planis inclinatis recepta} demonstrari posse,}}{%
\lemma{\textls{inclinatis}}%
\Bfootnote{%
\textit{(1)}~ex
\textit{(2)}~statuit objecto
\textit{(3)}~\textls{demonstrari}
\textit{(a)}~posse statuit objector,
\textit{(b)}~statuit objector posse, \textls{ex principio altiore}
\textit{(c)}~\textls{posse,}~\textit{L}}}
ex \textls{principio altiore},\indexS{principium altius} statuit 
\edtext{\indexN{\textls{Leibniz} (Leibnitius, GGL), Gottfried Wilhelm 1646\textendash1716}objector, nempe:}{%
\lemma{objector,}%
\Bfootnote{%
\textit{(1)}~quod
\textit{(2)}~nempe \textls{nullus} \bars{\textls{fiat} \textit{str. Hrsg.}} de
\textit{(3)}~nempe:~\textit{L}}}
\textls{quod nullus fiat descensus,\indexS{descensus} nisi plus descendatur,\indexS{descensus} quam ascenditur};\indexS{ascensus}
quod ex terminis patet; agitur enim descendendo, resistitur ascendendo,
itaque nisi plus agatur, quam resistitur, nihil agetur.
Non difficulter autem ostendi judicat: gravibus\indexS{grave} existentibus in situ et proportione 
secundum receptam legem\indexS{lex aequilibrii recepta} praescriptis, tunc sive in unam, 
\ledpb%
sive in alteram partem
descensum\indexS{descensus} fieri fingas, repertum iri, tantum esse descensum,\indexS{descensus} quantum ascensus,\indexS{ascensus}
adeoque nihil actum fuisse.
\pend
\pstart
Sed nos controversiae non misti virorum eruditorum\indexS{viri eruditi} judicio rem committimus.
\pend
